Point identification of the entire response vector means that the map
\[
 \widetilde{\mathcal M}\longmapsto\rho_x(\widetilde{\mathcal M}),
 \qquad
 \widetilde{\mathcal M}\in\mathfrak F_x^{\mathrm{law}}(P_X),
\]
is constant. By the sharp-effect-frontier theorem, its image is the nonempty set
\[
 \mathcal V_x(C)
 =\left\{\phi_x(j):j\in J_x(C)\right\},
 \qquad
 \phi_x(j)=\left(\frac{C_{ij}}{C_{xj}}\right)_{i\in\mathsf O\setminus\{x\}}. \tag{1}
\]
Thus the response vector is point identified if and only if \(|\mathcal V_x(C)|=1\). It remains to express this test directly in terms of \(J_x(C)\).

For every \(j\in J_x(C)\), the definition of \(J_x(C)\) gives \(C_{xj}\ne0\), so every division in (1) is valid. Suppose \(j,k\in J_x(C)\) and \(\phi_x(j)=\phi_x(k)\). Coordinatewise equality in (1) gives
\[
 \frac{C_{ij}}{C_{xj}}=\frac{C_{ik}}{C_{xk}}
 \quad\text{for every }i\in\mathsf O\setminus\{x\}. \tag{2}
\]
At coordinate \(x\), both normalized columns have value one. Adjoining this coordinate to (2) yields
\[
 \frac{c_j}{C_{xj}}=\frac{c_k}{C_{xk}},
 \qquad\text{hence}\qquad
 c_j=\frac{C_{xj}}{C_{xk}}c_k. \tag{3}
\]
The scalar in (3) is nonzero because both denominators and numerators are nonzero. Equation (3) says that \(c_j\) and \(c_k\) are proportional. The pairwise nonproportionality assumption therefore forces \(j=k\). Hence \(\phi_x\) is injective on the finite set \(J_x(C)\), and (1) implies
\[
 |\mathcal V_x(C)|=|J_x(C)|. \tag{4}
\]
Because the sharp-effect-frontier theorem also gives \(J_x(C)\ne\varnothing\), equations (1)--(4) prove the claimed vector criterion.

If that criterion fails, nonemptiness and (4) give distinct \(j,k\in J_x(C)\) with \(\phi_x(j)\ne\phi_x(k)\). The canonical-exact-law-witness lemma applies separately to these two indices. It gives canonical embeddings
\[
 \widetilde{\mathcal M}_j=(n,Q_j,H_j,\varepsilon),
 \qquad
 \widetilde{\mathcal M}_k=(n,Q_k,H_k,\varepsilon)
\]
in \(\mathfrak F_x^{\mathrm{law}}(P_X)\) and gives
\[
 \rho_x(\widetilde{\mathcal M}_j)=\phi_x(j)\ne
 \phi_x(k)=\rho_x(\widetilde{\mathcal M}_k). \tag{5}
\]
Membership in \(\mathfrak F_x^{\mathrm{law}}(P_X)\), by the definition of that fiber, proves separately that each system has observational law \(P_X\), that its structural matrix is invertible, and that its postintervention minor at \(x\) is invertible. Thus (5) supplies two exact same-observational-law, postintervention-solvable SCM countermodels; no stability restriction has been invoked.

We next prove the scalar corollary without replacing the jointly attainable vector set by a Cartesian product. The scalar part of the sharp-effect-frontier theorem gives
\[
 \left\{[\rho_x(\widetilde{\mathcal M})]_y:
 \widetilde{\mathcal M}\in\mathfrak F_x^{\mathrm{law}}(P_X)\right\}
 =\mathcal R_{xy}(C)
 =\left\{\frac{C_{yj}}{C_{xj}}:j\in J_x(C)\right\}. \tag{6}
\]
Consequently the \(y\)-coordinate is point identified if and only if \(|\mathcal R_{xy}(C)|=1\). Consider the two-row matrix \(C_{\{x,y\},J_x(C)}\). Since \(J_x(C)\ne\varnothing\) and \(C_{xj}\ne0\) for every \(j\in J_x(C)\), its \(x\)-row is nonzero and its rank is at least one. Its rank is at most one if and only if all its two-column minors vanish, namely if and only if
\[
 C_{xj}C_{yk}-C_{xk}C_{yj}=0
 \quad\text{for every }j,k\in J_x(C). \tag{7}
\]
For every such pair, division of (7) by the nonzero product \(C_{xj}C_{xk}\) is legitimate and gives
\[
 \frac{C_{yj}}{C_{xj}}=\frac{C_{yk}}{C_{xk}}. \tag{8}
\]
Thus (7) is equivalent to all elements of the nonempty set in (6) being equal. Combining the rank lower bound with (7)--(8) proves
\[
 |\mathcal R_{xy}(C)|=1
 \quad\Longleftrightarrow\quad
 \operatorname{rank}\!\left(C_{\{x,y\},J_x(C)}\right)=1. \tag{9}
\]
If (9) fails, (7) fails for some \(j,k\in J_x(C)\), so their ratios in (8) differ. Applying the canonical-exact-law-witness lemma to this pair gives two systems whose \(y\)-coordinate responses differ. Equation (1), however, shows that every whole response vector uses one common source index \(j\); equations (6)--(9) only take coordinate projections and never authorize cross-index recombination.

Finally, the witness algorithm shares one row-basis extension and one \(n\)-by-\(n\) matrix inverse. After those shared operations, each selected index requires at most one rank-one row shear, one perfect matching in an \((n-1)\)-by-\((n-1)\) nonzero pattern, and row permutation and normalization. A direct augmenting-path implementation performs at most \(n-1\) augmentations, each examining \(O(n^2)\) incidences, so one matching costs \(O(n^3)\) arithmetic and zero-test operations; the shear and normalization cost \(O(n^2)\). Repeating this for two fixed indices remains \(O(n^3)\), which proves the stated additional-operation bound.